% =====================================================================
\section{The cell in five minutes}
A cell is a droplet of aqueous solution inside a shell (the \emph{membrane}). Almost all the work inside it is done by \emph{proteins}. Some accelerate chemical reactions (these are called \emph{enzymes}), some serve as gates in the membrane, and some receive signals from outside (\emph{receptors}).

A protein is built from a gene in two steps: a working copy of the gene is made (mRNA), and the protein is assembled from it. The whole process is called gene \emph{expression}. A finished protein lives for a limited time and is then degraded. The amount of any protein is therefore a balance of synthesis and decay, and the cell governs itself by changing the rate of synthesis: it switches genes on and off in response to signals.

Energy is obtained in the \emph{mitochondria}, where nutrients are oxidised by oxygen. The energy is stored in ATP molecules and spent mainly on maintaining ion concentration differences across the membrane.
\facts
\begin{itemize}
\item Protein lifetimes range from minutes to weeks.\tU
\item The rate of an enzymatic reaction grows with the concentration of the substrate and saturates: there is a limited amount of enzyme, and each molecule of it is occupied for a finite time.\tU
\item The signal from a receptor reaches a gene along a chain of intermediate molecules; hours pass between the signal and an appreciable change in the amount of a protein.\tU
\item Most quantities in a cell are held in place by negative feedback.\tU
\item The brain is about 2\,\% of body mass and consumes about 20\,\% of the oxygen.\tU
\end{itemize}
\begin{ineq}
The amount of a protein $p$ obeys $p' = s(t) - p/\tau$: a first-order element whose time constant is the protein's lifetime. An enzymatic reaction is a saturating nonlinearity of the form $v = V s/(K + s)$. Control of genes supplies the slow variables, the chemistry of small molecules the fast ones. The saturation of the defence is exactly the source of the nonlinearity on which section~9 is built.
\end{ineq}
